\documentclass[12pt]{article}
\usepackage[utf8]{inputenc}
\usepackage[T1]{fontenc}
\usepackage{hyperref}
\title{Ufim Language Resources for K-12 Curriculum Development}
\author{Compiled by Claude Sonnet 5 (claude-sonnet-5) for Andr\'as Kornai}
\date{}
\begin{document}
\maketitle

\section{Introduction}
Ufim (ISO 639-3 ufi; Glottocode ufim1240) is a Finisterre language of Morobe Province, Papua New Guinea. Wikipedia cites about 1,300 speakers from a 2000 figure, in five villages north of Kaiapit, and Joshua Project gives 1,900 people. Many speakers also use Tok Pisin, and younger people increasingly rely on it. The searches found very little, so this report is short.

\section{Primary Educational Resources}
I found no primers, textbooks or school programmes for Ufim.

\section{Linguistic and Reference Materials}
\begin{itemize}
\item The Wikipedia article on Ufim summarises the available information. \url{https://en.wikipedia.org/wiki/Ufim_language}
\item The Max Planck Institute's numerals page for Ufim gives numerals from a 1988 linguistic study. \url{https://lingweb.eva.mpg.de/channumerals/Ufim.htm}
\end{itemize}

\section{Religious and Community Resources}
\begin{itemize}
\item Joshua Project has a profile of the Ufim people. It says Scripture translation has been started but is incomplete or hard to obtain. \url{https://joshuaproject.net/people_groups/15688/PP}
\end{itemize}

\section{Digital and Audio Resources}
I found none.

\section{Curriculum Development Context}
Almost nothing is documented. A curriculum would need new documentation work and an orthography, and would probably use Tok Pisin or English as a bridge language.

\end{document}